% ch06-ai-build
\section{الذكاء الاصطناعيّ (2): البناء والتشغيل}

في الفصل السابق فهمنا \emph{المفاهيم}: كيف يحوّل \en{Embedding} أي نصّ إلى متّجه
أرقام، وكيف يقيس \en{Cosine} قربَ متّجهين، وما معنى \en{RAG}، وكيف يوجّه النظام
الأسئلةَ بين ميزاته الثلاث (الشات، التوصيات، الأهليّة). في هذا الفصل ننزل إلى
\emph{المطبخ}: كيف بُنِي هذا كلّه فعليًّا في الكود، صفًّا صفًّا. هنمشي واحدة واحدة —
من بناء قاعدة المعرفة، إلى تخزين المتّجهات، إلى خطّ الشات الكامل، ثمّ التوصيات
والأهليّة الحسابيّتين، وأخيرًا شبكة الأمان والتكلفة.

\begin{idea}[الفكرة: ثلاث ميزات، محرّكان]
عندنا ثلاث ميزات ذكاء اصطناعيّ، لكن محرّكين فقط. \textbf{الشات} وحده يستخدم
\en{RAG} حقيقيًّا مع \en{GPT-4o-mini}. أمّا \textbf{التوصيات} و\textbf{الأهليّة}
فتعملان بخوارزميّة حسابيّة حتميّة (\en{deterministic}) بلا أيّ نداء شبكة ولا
\en{GPT} إطلاقًا. تذكّر هذا التقسيم؛ سيفسّر لك كثيرًا ممّا يأتي.
\end{idea}

% ═══════════════════════════════════════════════════════════════════════════
\subsection{قاعدة المعرفة: من خمسة مصادر إلى صفوف قابلة للبحث}

قبل أن يجيب الشات عن أيّ سؤال، لازم يكون عنده «كتاب» يقرأ منه. هذا الكتاب اسمه
\emph{قاعدة المعرفة} (\en{Knowledge Base})، ويبنيه صفٌّ واحد اسمه
\code{KnowledgeBaseIndexer}. مهمّته: يجمع البيانات من خمسة مصادر مختلفة، ويحوّل كلّ
عنصر إلى صفٍّ من نوع \code{KnowledgeDocument} في قاعدة \en{SQL Server}.

\begin{center}
\begin{tabular}{L{4.2cm} L{8.4cm}}
\toprule
\textbf{\color{spPrimary}المصدر} & \textbf{\color{spPrimary}ماذا يُفهرَس ومن أين} \\
\midrule
المنح المفتوحة & كلّ منحة حالتها \code{Open} — استعلام حيّ من القاعدة \\
الأسئلة الشائعة (\en{FAQ}) & مجموعة بيانات ثابتة من \code{KnowledgeBodies.cs} \\
الاستشاريّون الموثّقون & كلّ استشاريّ مُوثَّق وغير موقوف — استعلام حيّ \\
الموارد المنشورة & كلّ مورد حالته \code{Published} وغير محذوف \\
أفضل منشورات المجتمع & المواضيع التي صافي تقييمها 3 فأكثر (حتّى 200 موضوع) \\
\bottomrule
\end{tabular}
\end{center}

\begin{note}[ملاحظة: منشورات المجتمع مُطفأة افتراضيًّا]
المصدر الخامس (منشورات المجتمع) \textbf{لا} يدخل القاعدة إلّا لو فُعِّل المفتاح
\code{Ai:IncludeCommunityPostsInKb} صراحةً. السبب: هذه محتوى مستخدمين غير مُدقَّق،
ولا نريد للمساعد أن يُقدّم كلامَ أشخاص عشوائيّين وكأنّه معلومة رسميّة موثّقة، قبل أن
توجد آليّة مراجعة. لهذا عدد المستندات الفعليّ في الإنتاج يدور حول 692 مستندًا.
\end{note}

\subsubsection{الخوارزميّة: تحديث ذكيّ لا إعادة بناء كاملة}

الدالة \code{RebuildAsync} لا تمسح كلّ شيء وتبدأ من الصفر. بدلًا من ذلك تمشي في
أربع خطوات ذكيّة:

\begin{enumerate}
  \item \textbf{الإدراج والتحديث (\en{upsert})}: لكلّ مستند مطلوب، تحسب بصمة
    \en{SHA-256} على محتواه (\en{ContentEn + ContentAr}). لو المستند جديد تُدرِجه؛
    لو موجود تُحدّثه.
  \item \textbf{إزالة الأيتام (\en{orphans})}: أيّ مستند في القاعدة اختفى مصدره
    الأصليّ (مثلًا منحة أُغلقت) يُحذَف.
  \item \textbf{إعادة تضمين انتقائيّة}: المستند يُعاد تضمينه \emph{فقط} لو تغيّرت
    بصمته، أو كان تضمينه فارغًا، أو أُنشئ بنموذج مختلف. على قاعدة «دافئة» هذا صفرُ عمل.
  \item \textbf{التضمين على دُفعات}: المستندات المُعلَّقة تُرسَل إلى \en{Azure} في
    دُفعات من 16 مستندًا في المرّة.
\end{enumerate}

\begin{keyp}[نقطة مهمة: بصمة المحتوى هي مفتاح الرخص]
البصمة \en{SHA-256} هي حجر الزاوية. لو غيّرت وصف منحة واحدة، بصمتها فقط هي التي
تتغيّر، فيُعاد تضمين مستند واحد بتكلفة تكاد لا تُذكَر. أمّا الـ 691 مستندًا الباقون
فبصماتهم كما هي، فيُتخطَّون تمامًا. لهذا تكون عمليّة إعادة البناء الليليّة شبه مجّانيّة.
\end{keyp}

الآن نشوف الكود الحقيقيّ. أوّل مقتطف: خطوة الـ \en{upsert} مع فحص البصمة. لاحظ كيف
أنّ تغيّر البصمة \emph{يُفرِّغ} التضمين القديم حتّى يُعاد حسابه لاحقًا:

\begin{english}
\begin{lstlisting}[language={[Sharp]C}]
var hash = Sha256(d.EmbedText);              // SHA-256(ContentEn + ContentAr)
if (byKey.TryGetValue((d.SourceType, d.SourceKey), out var doc))
{
    doc.ContentEn = d.ContentEn;
    doc.ContentAr = d.ContentAr;
    // ...
    if (doc.ContentHash != hash)             // content actually changed?
    {
        doc.ContentHash = hash;
        doc.Embedding = [];                  // drop the stale vector
        doc.EmbeddingDimensions = 0;
        doc.EmbeddingModel = null;           // -> flagged for re-embedding
    }
}
\end{lstlisting}
\end{english}

شرح سطرًا بسطر: السطر الأوّل يحسب البصمة الجديدة. لو المستند موجود مسبقًا نُحدّث
حقوله. الشرط \code{doc.ContentHash != hash} يعني «هل تغيّر المحتوى حقًّا؟» — لو نعم،
نُفرِّغ حقل \code{Embedding} ونُصفّر النموذج، فيصير المستند «مُعلَّقًا» بانتظار إعادة
التضمين. لو المحتوى لم يتغيّر، لا نلمس التضمين، فلا تكلفة.

المقتطف الثاني: حلقة التضمين على دُفعات من 16. هنا يحدث الاتّصال الوحيد المدفوع مع
\en{Azure}، ثمّ تُحزَم النتيجة عبر \code{VectorMath.Pack}:

\begin{english}
\begin{lstlisting}[language={[Sharp]C}]
private const int EmbedBatchSize = 16;
// ...
for (var i = 0; i < pending.Count; i += EmbedBatchSize)
{
    var batch = pending.Skip(i).Take(EmbedBatchSize).ToList();
    var texts = batch.Select(EmbedText).ToList();

    var vectors = await embeddings.EmbedBatchAsync(texts, ct);   // Azure REST

    for (var j = 0; j < batch.Count && j < vectors.Count; j++)
    {
        batch[j].Embedding = VectorMath.Pack(vectors[j]);  // float[1536] -> byte[6144]
        batch[j].EmbeddingModel = model;
    }
}
\end{lstlisting}
\end{english}

الحلقة تتقدّم 16 عنصرًا في كلّ دورة (\code{i += EmbedBatchSize}). نأخذ دفعةً،
نُجهّز نصوصها، ونرسلها مرّةً واحدة إلى \code{EmbedBatchAsync}. الردّ متّجهات
\en{float}، ونحزم كلّ متّجه في عمود \en{VARBINARY} ونسجّل اسم النموذج الذي أنتجه.
إرسال 16 معًا أوفر بكثير من 16 نداءً منفصلًا.

\begin{secnote}[زاوية الأمان: بوّابة تمنع الصرف المزدوج]
عمليّة إعادة البناء تقرأ كلّ المستندات وتُضمِّن المُعلَّق منها مقابل مزوّد
\emph{مدفوع}. لو تشغّلت عمليّتان معًا (مثلًا تهيئة بدء التشغيل + طلب أدمن)، لضُمِّنت
المستنداتُ مرّتين — أي صرفٌ مضاعف. لهذا يحرس الكودُ العمليّةَ بقُفل ثابت على مستوى
العمليّة (\code{SemaphoreSlim RebuildGate}) يُسلسِل الطلبات فلا تتداخل أبدًا.
\end{secnote}

% ═══════════════════════════════════════════════════════════════════════════
\subsection{تخزين المتّجهات في \en{SQL Server} لا في قاعدة متّجهات}

قرارٌ معماريّ مهمّ: لم نستخدم قاعدة متّجهات متخصّصة (مثل \en{Pinecone} أو
\en{Qdrant})، بل خزّنّا المتّجهات داخل \en{SQL Server} نفسها. لكن \en{SQL} لا يعرف
نوع «متّجه»؛ يعرف \en{VARBINARY} فقط. فكيف نحوّل مصفوفة \en{float} إلى بايتات
والعكس؟ هذه مهمّة الصفّ \code{VectorMath}.

\begin{idea}[الفكرة: حسبة البايتات]
كلّ متّجه فيه 1536 رقمًا من نوع \en{float}. كلّ \en{float} يحتلّ 4 بايتات. إذن كلّ
مستند يحتاج $1536 \times 4 = 6144$ بايتًا فقط في عمود \en{VARBINARY}. رقمٌ صغير جدًّا،
ولهذا لا نحتاج بنيةً خاصّة لتخزينه.
\end{idea}

المقتطف التالي يُظهر التحويل في الاتّجاهين. لاحظ \code{Buffer.BlockCopy} — نسخٌ خامّ
للذاكرة بلا حلقات، سريعٌ للغاية:

\begin{english}
\begin{lstlisting}[language={[Sharp]C}]
public static byte[] Pack(float[] vector)
{
    if (vector.Length == 0) return [];
    var bytes = new byte[vector.Length * sizeof(float)];   // 1536 * 4 = 6144
    Buffer.BlockCopy(vector, 0, bytes, 0, bytes.Length);
    return bytes;
}

public static float[] Unpack(byte[] bytes)
{
    if (bytes.Length == 0) return [];
    var vector = new float[bytes.Length / sizeof(float)];
    Buffer.BlockCopy(bytes, 0, vector, 0, vector.Length * sizeof(float));
    return vector;
}
\end{lstlisting}
\end{english}

الدالة \code{Pack} تحجز مصفوفة بايتات بطول $\text{count} \times 4$، ثمّ تنسخ
الذاكرة نسخًا مباشرًا. \code{Unpack} تفعل العكس تمامًا: تقسم عدد البايتات على 4 لتعرف
عدد الأرقام، ثمّ تنسخها إلى مصفوفة \en{float} جاهزة لحساب \en{Cosine}.

\begin{why}
عند 692 مستندًا، مسحُ القاعدة كلّها في الذاكرة وحسابُ \en{Cosine} لكلّ مستند يكتمل في
أقلّ من 10 مِلّي ثانية. أيّ قاعدة متّجهات خارجيّة كانت ستضيف خدمةً جديدة تُنشَر
وتُدار، وقفزةَ شبكة إضافيّة، وتكلفةً تشغيليّة — كلّ ذلك لِيصير أبطأ لا أسرع عند هذا
الحجم. عتبةُ الانتقال إلى قاعدة متّجهات حقيقيّة تبدأ عند عشرات الآلاف من المستندات.
\end{why}

% ═══════════════════════════════════════════════════════════════════════════
\subsection{خطّ الشات الكامل: تسع خطوات}

الآن الجزء الأهمّ: ماذا يحدث بالضبط عندما يكتب الطالب سؤالًا في الشات؟ الرحلة تمرّ
بتسع خطوات مرتّبة. الجدول التالي يلخّصها، ثمّ نشرح المحطّات الحسّاسة:

\begin{center}
\begin{tabular}{C{0.7cm} L{4.4cm} L{7.2cm}}
\toprule
\textbf{\color{spPrimary}\#} & \textbf{\color{spPrimary}الخطوة} & \textbf{\color{spPrimary}ماذا يحدث} \\
\midrule
1 & استقبال الرسالة & الطالب يرسل سؤاله النصّيّ \\
2 & كشف اللغة & \code{RagSupport.IsArabic()} يقرّر: عربيّ أم إنجليزيّ \\
3 & فحص الإعداد & هل \en{Azure} مُهيّأ؟ لو لا $\rightarrow$ سقوط محلّيّ فورًا \\
4 & تضمين السؤال & \code{EmbedAsync()} يحوّل السؤال إلى \en{float[1536]} \\
5 & الاسترجاع & حساب \en{Cosine} على 692 مستندًا، وأخذ أعلى 5 \\
6 & التصفية & إبقاء المستندات التي درجتها 0.6 فأكثر فقط \\
7 & بناء السياق & \code{BuildContextBlock()} ينسّق المستندات المختارة \\
8 & نداء \en{GPT} & إرسال: التعليمات + التاريخ + السياق + السؤال \\
9 & بناء الاستجابة & نصّ الجواب + قائمة المصادر + التكلفة المقدّرة \\
\bottomrule
\end{tabular}
\end{center}

الخطوات 4 و5 و6 هي قلب \en{RAG}. المقتطف التالي — من صفّ الاسترجاع
\code{KnowledgeRetriever} — يُظهر كيف نحمّل المستندات المُضمَّنة \emph{بالنموذج
النشِط فقط}، ثمّ نرتّبها بالدرجة ونأخذ أعلى \en{topK}:

\begin{english}
\begin{lstlisting}[language={[Sharp]C}]
var docs = await db.KnowledgeDocuments.AsNoTracking()
    .Where(d => d.EmbeddingModel == model && d.EmbeddingDimensions > 0)
    .Select(d => new { d.Id, d.TitleEn, d.ContentEn, d.Embedding /* ... */ })
    .ToListAsync(ct);

var queryVector = await embeddings.EmbedAsync(query, ct);

return [.. docs
    .Select(d => new {
        Doc = d,
        Score = VectorMath.CosineSimilarity(queryVector, VectorMath.Unpack(d.Embedding)),
    })
    .OrderByDescending(x => x.Score)
    .Take(topK)];                            // top-5 by default
\end{lstlisting}
\end{english}

الشرط \code{EmbeddingModel == model} مهمّ: متّجهاتُ نموذجٍ مختلف تعيش في «فضاء» آخر
ولا يصحّ مقارنتها، فنتجاهلها. نحسب \en{Cosine} بين متّجه السؤال وكلّ مستند بعد فكّ
حزمته بـ \code{Unpack}، نرتّب تنازليًّا، ونأخذ الأعلى. القيمة الافتراضيّة لـ
\code{Ai:RagTopK} هي 5، ولـ \code{Ai:RagMinScore} هي 0.6.

\subsubsection{تصميم تعليمات النظام (\en{System Prompt})}

قبل نداء \en{GPT}، نُرسِل «تعليمات نظام» تُشكّل شخصيّة المساعد (بنسختين: عربيّة
وإنجليزيّة). أهمّها أن يُطلَب منه:

\begin{itemize}
  \item أن يذكر المنحَ والاستشاريّين والموارد \textbf{بالاسم} حين تَرِد في كتلة السياق.
  \item أن يساعد في صياغة خطاب الغرض وخطابات التوصية والسيرة الذاتيّة عند الطلب.
  \item ألّا \textbf{يخترع} أبدًا اسمَ منحةٍ أو موعدًا أو مبلغًا غير موجود في السياق.
  \item ألّا \textbf{يرفض} طلبًا معقولًا من المستخدم.
\end{itemize}

\subsubsection{ذاكرة الجلسة (\en{Session Memory})}

كيف «يتذكّر» الشات ما قيل قبل قليل بلا مخزنٍ على الخادم؟ ببساطة: قبل كلّ نداء نُضيف
أدوار المحادثة السابقة (آخر ~20 دورًا) بين تعليمات النظام والسؤال الحاليّ:

\begin{english}
\begin{lstlisting}[language={[Sharp]C}]
var messages = new List<object>
{
    new { role = "system", content = system },
};
foreach (var turn in history)                     // prior turns this session
    messages.Add(new { role = turn.Role, content = turn.Content });
messages.Add(new { role = "user", content = user });   // current question
\end{lstlisting}
\end{english}

القائمة تبدأ بتعليمات النظام، ثمّ نُضيف كلّ دورٍ سابق بدوره (\code{user} أو
\code{assistant})، ثمّ نُلحِق السؤالَ الحاليّ. النموذج يقرأ السياق كاملًا فيبدو
«متذكّرًا». نداء \en{GPT-4o-mini} يمرّ بـ \code{temperature=0.3} — قيمة منخفضة تجعل
الإجابات ثابتةً ومنضبطة لا خياليّة.

% ═══════════════════════════════════════════════════════════════════════════
\subsection{التوصيات: درجاتٌ حسابيّة بلا \en{GPT}}

هنا ننتقل إلى المحرّك الثاني. خلافًا لما قد يظنّه كثيرون، توليدُ التوصيات
\textbf{ليس} فيه أيّ ذكاء اصطناعيّ توليديّ. إنّه صفٌّ حسابيّ اسمه
\code{LocalAiService} يُعطي كلّ منحةٍ مفتوحة \emph{درجةً} من 100 بناءً على مطابقتها
لملفّ الطالب.

\begin{secnote}[زاوية الأمان: لا شبكة، لا \en{GPT}، لا مفاجآت]
لا يوجد نداء \en{GPT}، ولا \en{RAG}، ولا أيّ طلب \en{HTTP} في خطّ التوصيات إطلاقًا.
كلّه يعمل داخل الخادم بشكل حتميّ. النتيجة نفسها للمدخلات نفسها دائمًا — قابلة للاختبار
ولا تكلّف شيئًا ولا تتعطّل بسبب مزوّد خارجيّ.
\end{secnote}

معادلة الدرجة تجمع أربعة مكوّنات بأوزان مختلفة:

\[
\text{Score} \;=\; \underbrace{S_{\text{level}}}_{\le 30}
\;+\; \underbrace{S_{\text{field}}}_{\le 40}
\;+\; \underbrace{S_{\text{country}}}_{\le 25}
\;+\; \underbrace{B_{\text{funding}}}_{\le 5}
\;\in\; [0,\,100]
\]

مجال الدراسة هو الوزن الأكبر (حتّى 40 نقطة) لأنّ توافق التخصّص أهمّ من الجغرافيا.
المقتطف الحقيقيّ من دالة التوليد:

\begin{english}
\begin{lstlisting}[language={[Sharp]C}]
var score = 0;
if (userLevel.HasValue && userLevel.Value == c.TargetLevel) score += 30;
score += OverlapScore(preferredFields, fosFields, maxPoints: 40);
// countries compared on a canonical ISO key ("Egypt"/"masr" -> "EG")
score += OverlapScore(NormalizeCountries(preferredCountries),
                      NormalizeCountries(countries), maxPoints: 25);
if (c.FundingAmountUsd >= 10_000) score += 5;      // funding bonus
scored.Add((c.Id, Math.Min(score, 100), c.TitleEn, c.TitleAr));
\end{lstlisting}
\end{english}

المستوى الأكاديميّ: تطابقٌ تامّ يعطي 30. المجال والدولة: عبر دالّة
\code{OverlapScore} التي تحسب نسبة التقاطع مضروبةً في السقف. مكافأةُ 5 نقاط لو كان
التمويل 10000 دولار فأكثر. أخيرًا \code{Math.Min(score, 100)} تضمن ألّا تتجاوز
الدرجةُ 100. ثمّ تُترجَم الدرجة إلى تسمية مفهومة للطالب:

\begin{center}
\begin{tabular}{C{2.6cm} L{3.0cm} L{6.6cm}}
\toprule
\textbf{\color{spPrimary}الدرجة} & \textbf{\color{spPrimary}التسمية} & \textbf{\color{spPrimary}ما يظهر للطالب} \\
\midrule
80--100 & توافق قويّ  & قريبة من اهتماماتك ومستواك الأكاديميّ \\
50--79  & توافق جزئيّ & تتقاطع مع بعض تفضيلاتك \\
20--49  & توافق محدود & تلمس مجالاتك، تستحقّ نظرةً سريعة \\
0--19   & توافق ضعيف  & لا تتماشى حاليًّا مع ملفّك \\
\bottomrule
\end{tabular}
\end{center}

\subsubsection{التخزين المؤقّت وإعادة الترطيب (\en{Re-hydration})}

حسابُ الدرجات يجري مرّةً وتُخزَّن النتيجة \emph{مؤقّتًا} كـ \en{JSON} داخل الحقل
\code{ResponseText} من كيان \code{AiInteraction}. لكن هنا حيلةٌ ذكيّة: الـ \en{JSON}
المخزَّن يحفظ فقط الحقولَ التي أنتجها الذكاء (المعرّف، الدرجة، الشرح). أمّا
\emph{الموعد النهائيّ} و\emph{مبلغ التمويل} فيُجلَبان \textbf{حيًّا} من القاعدة عند
كلّ قراءة.

\begin{keyp}[نقطة مهمة: لماذا نُعيد الترطيب؟]
لو خزّنّا الموعد النهائيّ داخل الكاش، ثمّ غيّره الأدمن، لظهر للطالب موعدٌ قديمٌ خاطئ.
بإعادة الترطيب من القاعدة حيًّا، الدرجة المحفوظة تبقى (فهي لا تتغيّر كثيرًا)، بينما
البياناتُ الحسّاسة للوقت تكون دائمًا حديثة. الكاش صالح 24 ساعة؛ بعدها تُولَّد توصياتٌ
جديدة.
\end{keyp}

% ═══════════════════════════════════════════════════════════════════════════
\subsection{الأهليّة: مطابقة الملفّ الشخصيّ}

المحرّك الحسابيّ نفسه يخدم الميزة الثالثة: فحص الأهليّة. هنا نقارن ملفّ الطالب
بمتطلّبات منحةٍ واحدة عبر ثلاثة معايير:

\begin{center}
\begin{tabular}{L{3.2cm} L{4.4cm} L{4.6cm}}
\toprule
\textbf{\color{spPrimary}المعيار} & \textbf{\color{spPrimary}مصدر الطالب} & \textbf{\color{spPrimary}منطق المطابقة} \\
\midrule
المستوى الأكاديميّ & \code{AcademicLevel} & تطابق تامّ \\
الدولة & الدول المفضّلة + الجنسيّة & عبر \code{CountryNormalizer} \\
مجال الدراسة & \code{FieldOfStudy} & احتواء غير حسّاس للحالة \\
\bottomrule
\end{tabular}
\end{center}

كلّ معيار يأخذ إحدى القيم: \code{yes} أو \code{partial} أو \code{no} أو
\code{unknown} (حين يكون حقلُ الطالب فارغًا). ثمّ تُجمَع هذه القيم في حكمٍ واحد عبر
دالّة \code{DeriveVerdict}:

\begin{english}
\begin{lstlisting}[language={[Sharp]C}]
internal static EligibilityVerdict DeriveVerdict(
    IReadOnlyList<AiEligibilityCriterion> criteria)
{
    if (criteria.Count == 0) return EligibilityVerdict.NotEligible;

    if (criteria.Any(c => c.Match == "no"))            // any hard fail
        return EligibilityVerdict.NotEligible;

    if (criteria.Any(c => c.Match == "partial" || c.Match == "unknown"))
        return EligibilityVerdict.PartiallyEligible;   // any uncertainty

    return EligibilityVerdict.Eligible;                // all clear
}
\end{lstlisting}
\end{english}

المنطق بسيطٌ ومتحفّظ: أيُّ معيار قيمته \code{no} كافٍ وحده ليكون الحكم
\code{NotEligible}. لو لا يوجد رفضٌ صريح لكن يوجد \code{partial} أو \code{unknown}،
فالحكم \code{PartiallyEligible}. ولا نقول \code{Eligible} إلّا لو كانت كلّ المعايير
\code{yes}.

\begin{note}[ملاحظة: تطبيع أسماء الدول ليس بديهيًّا]
منحةٌ تكتب دولتها «\en{France}» يجب أن تطابق طالبًا كتب «\en{FR}» أو «\en{france}»
أو «فرنسا». الدالّة \code{CountryNormalizer.Matches} تُرجِع كلّ الصيغ إلى مفتاحٍ
معياريّ واحد قبل المقارنة. بدون هذا، لقرأ طالبٌ مصريّ يفضّل ألمانيا كلمة «\en{no}»
على منحةٍ كتبت الدولة «\en{DE}».
\end{note}

% ═══════════════════════════════════════════════════════════════════════════
\subsection{شبكة الأمان: السقوط الآمن (\en{Fallback}) ثلاثيّ المستوى}

ماذا لو تعطّل \en{Azure} تمامًا؟ لا يجوز أن يرى المستخدمُ خطأً. لهذا كلّ نداء شات
مُحاطٌ بسلسلة سقوطٍ من ثلاثة مستويات، أيُّ فشلٍ فيها يوجّه الطلبَ إلى
\code{LocalAiService.AskAsync}:

\begin{center}
\begin{tabular}{C{0.7cm} L{5.4cm} L{6.2cm}}
\toprule
\textbf{\color{spPrimary}\#} & \textbf{\color{spPrimary}نقطة الفحص} & \textbf{\color{spPrimary}عند الفشل} \\
\midrule
1 & هل النقطة والمفتاح مُهيّآن؟ & سقوطٌ محلّيّ فوريّ \\
2 & هل نجح نداء الاسترجاع؟ & سقوطٌ محلّيّ (خطأ \en{HTTP}) \\
3 & هل نجح نداء \en{GPT}؟ & سقوطٌ محلّيّ (خطأ أو مهلة) \\
\bottomrule
\end{tabular}
\end{center}

والجميل أنّ المستوى المحلّيّ ليس ردًّا أخرقَ ثابتًا. \code{AskAsync} المحلّيّة تُجرِّب
\emph{مسارين}: أوّلًا نفس استرجاع \en{RAG} (بنفس \code{KnowledgeRetriever}) لكن بجواب
\emph{استخلاصيّ} — تعرض أقربَ مستند مفهرَس بدل توليد نصّ. ولو لم تجد شيئًا فوق العتبة
0.6، تسقط إلى موجّهِ كلماتٍ مفتاحيّة (\en{keyword router}) يفهم نيّة السؤال (موعد،
أهليّة، توصية...) ويردّ ردًّا مفيدًا. هكذا يبقى \en{RAG} قابلًا للعرض حتّى بلا أيّ
مفتاح \en{API}.

% ═══════════════════════════════════════════════════════════════════════════
\subsection{التكلفة والضبط الدقيق وكيان \code{AiInteraction}}

\subsubsection{نموذج التكلفة وسقف الإنفاق}

نموذج \en{GPT-4o-mini} رخيصٌ جدًّا. التكلفة تُحسَب من عدد الرموز (\en{tokens}) داخلًا
وخارجًا:

\[
\text{Cost} \;=\; N_{\text{prompt}} \times \frac{0.15}{10^6}
\;+\; N_{\text{completion}} \times \frac{0.60}{10^6}
\]

(أي 0.15 دولار لكلّ مليون رمز دخْل، و0.60 دولار لكلّ مليون رمز خرْج). الجدول يوضّح
حجم الأرقام العمليّ:

\begin{center}
\begin{tabular}{L{5.0cm} C{2.2cm} C{2.2cm} C{2.6cm}}
\toprule
\textbf{\color{spPrimary}السيناريو} & \textbf{الدخل} & \textbf{الخرج} & \textbf{التكلفة} \\
\midrule
دورة شات بسيطة & 300 & 150 & \$0.000135 \\
سياق غنيّ (5 مستندات) & 600 & 300 & \$0.000270 \\
1000 دورة يوميًّا & --- & --- & \$0.27 تقريبًا \\
سقف \code{DailyCostLimitUsd} & \multicolumn{3}{L{7.0cm}}{\$5.00 --- حدّ يوميّ صارم في \code{AiOptions}} \\
\bottomrule
\end{tabular}
\end{center}

\begin{why}
حتّى لو خرجت الأمور عن السيطرة، المفتاح \code{Ai:DailyCostLimitUsd} يضع سقفًا صلبًا
عند خمسة دولارات في اليوم. عند بلوغه يتوقّف النظام عن نداء \en{Azure} ويسقط تلقائيًّا
إلى المحرّك المحلّيّ المجّانيّ. هكذا لا يمكن لخطأٍ أو إساءة استخدام أن يُنتِج فاتورةً
مفاجئة.
\end{why}

\subsubsection{مسار الضبط الدقيق (\en{Fine-Tuning}) داخل لوحة الأدمن}

جُهِّزت مجموعةُ بياناتٍ للضبط الدقيق (نحو 332 مثالًا بصيغة محادثات \en{OpenAI}).
والمسار كلّه يُدار من لوحة الأدمن في أربع خطوات: \emph{تصدير} الـ \en{JSONL}، ثمّ
\emph{رفعه}، ثمّ \emph{تدريب} المهمّة على \en{Azure}، ثمّ \emph{تفعيل} النموذج
الناتج. وأجمل ما فيه أنّ تفعيله لا يحتاج تغييرَ كود؛ يكفي أن يختار الكودُ اسمَ
النشرِ بأولويّةٍ ثلاثيّة:

\begin{english}
\begin{lstlisting}[language={[Sharp]C}]
// Priority: DB platform setting -> appsettings override -> base model
var deployment = !string.IsNullOrWhiteSpace(dbDeployment)
    ? dbDeployment                       // e.g. "scholarpath-ft-001"
    : !string.IsNullOrWhiteSpace(az.FineTunedDeploymentName)
        ? az.FineTunedDeploymentName
        : az.DeploymentName;             // "gpt-4o-mini"  (base)
\end{lstlisting}
\end{english}

أوّلًا يُقرأ اسمُ النشر من إعدادات المنصّة في القاعدة؛ فإن كان فارغًا يُجرَّب إعداد
الملفّ؛ وإلّا يُستخدَم النموذجُ الأساس \en{gpt-4o-mini}. الأدمن يستطيع «الرجوع للنموذج
الأساس» بضغطة واحدة بلا إعادة تشغيل.

\subsubsection{كيان \code{AiInteraction}: سجلٌّ وكاشٌ في جدولٍ واحد}

كلّ عمليّات الذكاء الاصطناعيّ (شات، توصية، أهليّة) تُسجَّل في جدولٍ واحد هو
\code{AiInteraction}. إنّه في آنٍ معًا \textbf{سجلّ تدقيق} و\textbf{طبقة تخزين مؤقّت}:

\begin{english}
\begin{lstlisting}[language={[Sharp]C}]
public class AiInteraction : AuditableEntity
{
    public Guid UserId { get; set; }
    public AiFeature Feature { get; set; }   // Recommendation | Chat | Eligibility
    public AiProvider Provider { get; set; } // AzureOpenAi | Local
    public string? SessionId { get; set; }   // groups chat turns into a session

    public string PromptText { get; set; }   // what was asked
    public string ResponseText { get; set; } // answer text OR cached JSON
    public string? MetadataJson { get; set; }// eligibility breakdown, etc.

    public int PromptTokens { get; set; }
    public int CompletionTokens { get; set; }
    public decimal CostUsd { get; set; }     // billing
}
\end{lstlisting}
\end{english}

المفتاح لفهم هذا الكيان أنّ معنى الحقل \code{ResponseText} \emph{يتغيّر} حسب
\code{Feature}: في الشات هو نصّ الجواب؛ في التوصيات هو الـ \en{JSON} المخزَّن مؤقّتًا؛
في الأهليّة هو الحكم، وتفاصيله في \code{MetadataJson}. جدولٌ واحد، ثلاثة استعمالات —
بسيطٌ وأنيق.

% ═══════════════════════════════════════════════════════════════════════════
\begin{summary}[الخلاصة في سطور]
\begin{itemize}
  \item \textbf{قاعدة المعرفة} تُبنى من خمسة مصادر (منح، أسئلة شائعة، استشاريّون،
    موارد، منشورات مجتمع) إلى صفوف \code{KnowledgeDocument}؛ بصمة \en{SHA-256} تجعل
    إعادة البناء انتقائيّة ورخيصة، والتضمين يجري في دُفعاتٍ من 16.
  \item \textbf{المتّجهات} تُخزَّن في \en{SQL Server} كـ \en{VARBINARY} بحجم
    $1536 \times 4 = 6144$ بايتًا لكلّ مستند؛ كافٍ تمامًا عند 692 مستندًا، والمسح
    يكتمل في أقلّ من 10 مِلّي ثانية بلا قاعدة متّجهات.
  \item \textbf{خطّ الشات} تسع خطوات: كشف لغة، فحص إعداد، تضمين، استرجاع \en{Cosine}،
    تصفية عند 0.6، بناء سياق، نداء \en{GPT-4o-mini} بـ \code{temperature=0.3}، ثمّ
    بناء الجواب والمصادر؛ مع ذاكرة جلسةٍ من آخر ~20 دورًا.
  \item \textbf{التوصيات والأهليّة} حسابيّتان بلا \en{GPT}: درجةٌ من
    (مستوى 30 + مجال 40 + دولة 25 + مكافأة 5)، وحكمُ أهليّةٍ من ثلاثة معايير مع تطبيع
    الدول؛ التوصيات تُخزَّن مؤقّتًا وتُعاد ترطيبُ موعدها وتمويلها حيًّا من القاعدة.
  \item \textbf{شبكة الأمان} ثلاثيّة المستوى تسقط دائمًا إلى \code{LocalAiService}؛
    وسقفُ \code{DailyCostLimitUsd} عند خمسة دولارات يمنع أيّ فاتورةٍ مفاجئة؛ وكلّ
    عمليّة تُسجَّل في كيان \code{AiInteraction} الذي يجمع بين السجلّ والكاش.
\end{itemize}
\end{summary}
